% Part: incompleteness
% Chapter: incompleteness-provability
% Section: godels-paper

\documentclass[../../../include/open-logic-section]{subfiles}

\begin{document}

\olfileid[fr]{inc}{inp}{gop}

\olsection{Comparaison avec l'article original de Gödel}

Il vaut la peine de consacrer un peu de temps à
l'article de Gödel de 1931. Son introduction esquisse
les idées que nous venons d'étudier. Même en le
parcourant rapidement, on distingue les étapes de
l'argument : Gödel décrit d'abord le système
formel~$P$, sa syntaxe, ses axiomes et ses règles
de preuve. Il définit ensuite les fonctions et les
relations récursives primitives, montre que
$x B y$ est récursive primitive et établit que ces
fonctions et relations sont représentables dans
$\Th{P}$. Il démontre alors le théorème
d'incomplétude comme ci-dessus. Dans la section~3,
il montre que l'énoncé indémontrable peut être un
!!{sentence} du langage de l'arithmétique. C'est
l'origine du lemme de la fonction~$\beta$, que
nous avons également utilisé pour traiter les
suites dans la preuve que les fonctions récursives totales
sont représentables dans~$\Th{Q}$. Gödel n'isole
pas un ensemble minimal d'axiomes suffisant pour
cette construction; nous savons maintenant que
$\Th{Q}$ suffit. Enfin, dans la section~4,
il esquisse une preuve du second théorème
d'incomplétude.

\end{document}
